\subsection{中央の二項係数との関係}

偶数公式の分数には，階乗を使った別の形がある。
\[
 (2m)!!=2^m m!,\qquad
 (2m-1)!!=\frac{(2m)!}{2^m m!}
\]
だから
\begin{equation}\label{式:ウォリスと二項係数}
 I_{2m}=\frac{\pi}{2}\frac{(2m)!}{2^{2m}(m!)^2}
 =\frac{\pi}{2}\frac{\binom{2m}{m}}{4^m}.
\end{equation}
ここで$\binom{2m}{m}=\frac{(2m)!}{(m!)^2}$は二項係数である。

さらに偶数公式と奇数公式を掛けると，多くの因子が消えて
\[
 I_{2m}I_{2m+1}=\frac{\pi}{2(2m+1)}
\]
となる。これと$\frac{I_{2m+1}}{I_{2m}}\to1$を組み合わせれば
\[
 mI_{2m}^2
 =\frac{m\pi}{2(2m+1)}\frac{I_{2m}}{I_{2m+1}}
 \longrightarrow\frac{\pi}{4}.
\]
$I_{2m}>0$なので平方根を取ることができ，
\[
 \sqrt m\,I_{2m}\longrightarrow\frac{\sqrt\pi}{2}.
\]
式\eqref{式:ウォリスと二項係数}に戻すと
\begin{equation}\label{式:中央二項係数の極限}
 \lim_{m\to\infty}\frac{\sqrt m}{4^m}\binom{2m}{m}
 =\frac1{\sqrt\pi}
\end{equation}
を得る。階乗を含む数列の大きさが，三角関数の積分を経由して分かったことになる。

